\documentclass[tikz,border=10pt]{standalone}
\input{../../docs/tikz/dag_preamble.tex}

\begin{document}
\begin{tikzpicture}[
  x=1cm,
  y=1cm,
  every node/.append style={outer sep=5pt, font=\sffamily\small}
]

\dagPaperMetadata{Designing Optimal Binary Rating Systems}{Garg, Johari}{AISTATS / PMLR 89}{2019}{https://proceedings.mlr.press/v89/garg19a.html}

\dagPaperLegendRightOfMetadata{
\node[dag_model, dag_template_legend] (legDef) at (0,0) {definition /\\formula};
\node[dag_result, dag_template_legend] (legDone) at (4.2,0) {formalized\\result};
}
\daglegend{(legDef)(legDone)}{Legend}

\begin{dagPaperBody}

\node[dag_model] (KL) at (0,0) {
  \textbf{Bernoulli KL (Sec. 3.3)} \\
  finite-support KL and support-safe convention
};

\node[dag_result] (AdjRate) at (6.2,0) {
  \textbf{Theorem 3.1 Rate Formula} \\
  adjacent binary threshold rate in real and support-safe forms
};

\node[dag_result] (ClosedRate) at (12.4,0) {
  \textbf{Lemma 3.1 Closed Rate} \\
  weighted common-threshold expression realizes the displayed value
};

\node[dag_result] (LevelOpt) at (18.6,0) {
  \textbf{Lemma 3.1 Finite Levels} \\
  endpoint-aware equalized levels exist, are unique, and are maximin optimal
};

\node[dag_result] (FiniteRate) at (0,-3.3) {
  \textbf{Finite Adjacent Objective} \\
  endpoint/source-shaped rows construct finite certificates and give exact rates
};

\node[dag_result] (T31Value) at (0,-5.25) {
  \textbf{Theorem 3.1 Value Logic} \\
  bounded convergence and two-stage value/rate lexicographic optimality
};

\node[dag_result] (C5) at (6.2,-3.3) {
  \textbf{Lemma C.5} \\
  doubled endpoint chain is feasible and equalizes adjacent rates
};

\node[dag_result] (C2C8) at (12.4,-3.3) {
  \textbf{Cor. C.2; Lemmas C.6-C.8} \\
  common rate goes to zero, mesh shrinks, and finite support bounds hold
};

\node[dag_result] (C9) at (18.6,-6.6) {
  \textbf{Lemma C.9} \\
  nested-bisection operation count satisfies the step bound
};

\node[dag_result] (T32Cert) at (18.6,-3.3) {
  \textbf{Theorem 3.2 Finite Certificate} \\
  calculated-grid early-return endpoint gives additive loss and runtime bounds
};

\node[dag_result] (TC1) at (0,-7.3) {
  \textbf{Theorem C.1 Weighted Skeleton} \\
  uniform, normalized-log, or pointwise local certificates give weighted integral rates
};

\node[dag_result] (C3C4) at (6.2,-7.3) {
  \textbf{C.3/C.4 Fixed Discretization} \\
  source-defined $Wbar_k$ and ordered rectangles have the adjacent exponent
};

\node[dag_result] (C3Finite) at (6.2,-9.45) {
  \textbf{C.3 Finite Algebra} \\
  endpoint-pair finite sums, adjacent dominance, and partition aggregation
};

\node[dag_result] (T31) at (12.4,-7.3) {
  \textbf{Full Theorem 3.1} \\
  continuum optimal discretization and maximin decomposition are checked
};

\node[dag_result] (C4Witness) at (0,-11.6) {
  \textbf{C.4 Source Bridge} \\
  monotone-source and positive-support primitives feed the reverse obstruction
};

\node[dag_result] (C4Reverse) at (6.2,-11.6) {
  \textbf{Lemma C.4 Reverse Bridge} \\
  local-constancy and finite-range tie-erased conventions rule out positive rates
};

\node[dag_result] (TB1) at (12.4,-10.9) {
  \textbf{Theorem B.1 Bridge} \\
  quantile-floor representation and uniform convergence give uniform limits
};

\node[dag_result] (C11) at (18.6,-8.75) {
  \textbf{Lemmas C.10-C.12} \\
  Spearman reduction and Kendall/Spearman finite objectives favor equispaced cutpoints
};

\node[dag_result] (Examples) at (18.6,-11.5) {
  \textbf{Definition C.1; Lemma C.10; Cor. C.4} \\
  finite reductions and equispaced Kendall/Spearman bridge are checked
};

\node[dag_template_note] (ScopeNote) at (9.3,-15.4) {
\textbf{Scope note.} Every box in this DAG is a Lean-checked definition,
  theorem, or result row. Some theorem boxes keep explicit grid, rate, weight,
  essential-infimum, local pointwise-rate, or iteration hypotheses visible in
  the paper-facing interface; these are closed theorem statements, not open
  assumptions or partial nodes.
  The reusable weighted Laplace skeleton behind
  Theorem C.1 and the endpoint-pair finite-sum, partition-integral,
  adjacent-dominance, source-defined fixed-discretization $Wbar_k$, and
  positive-rate algebra behind Lemmas C.3-C.4 and Theorem 3.1's
  fixed-discretization bridge are formalized under visible hypotheses.
  Corollary C.2, the finite Lemma C.10-C.12 example reductions, the reverse
  C.4 obstruction core, the concrete monotone-interval/local-constancy and
  finite-range pointwise-LDP reverse bridges, the Theorem 3.2
  calculated-grid early-return endpoint, and the common floor-coordinate /
  canonical B.1/Corollary C.4 convergence branches are formalized. The full
  Theorem 3.1 continuum discretization path and the Kendall/Spearman
  convergence examples are closed in the paper-facing interface.
  The finite Appendix B.2/B.3 learning wrappers are Lean-checked; the
  empirical/visualization material remains source scope but is not a
  theorem-DAG target here.
};

\draw[dag_arrow] (KL.east) -- (AdjRate.west);
\draw[dag_arrow] (AdjRate.east) -- (ClosedRate.west);
\draw[dag_arrow] (ClosedRate.east) -- (LevelOpt.west);

\draw[dag_arrow] (KL.south) -- (FiniteRate.north);
\draw[dag_arrow] (AdjRate.south west) to[out=-135,in=55,looseness=1.0] (FiniteRate.north east);

\draw[dag_arrow] (LevelOpt.south west) to[out=-120,in=35,looseness=1.05] (C5.north east);
\draw[dag_arrow] (C5.east) -- (C2C8.west);
\draw[dag_arrow] (C2C8.east) -- (T32Cert.west);
\draw[dag_arrow] (C9.north) -- (T32Cert.south);

\draw[dag_arrow] (TC1.east) -- (C3C4.west);
\draw[dag_arrow] (C3Finite.north) -- (C3C4.south);
\draw[dag_arrow] (C3C4.east) -- (T31.west);
\draw[dag_arrow] (C3C4.west) to[out=180,in=105,looseness=1.1] (C4Reverse.north west);
\draw[dag_arrow] (C4Witness.east) -- (C4Reverse.west);
\draw[dag_arrow] (C4Reverse.east) -- (T31.south west);
\draw[dag_arrow] (FiniteRate.south) -- (T31Value.north);
\draw[dag_arrow] (T31Value.east) to[out=0,in=165,looseness=1.0] (T31.west);
\draw[dag_arrow] (FiniteRate.south east) to[out=-45,in=145,looseness=1.05] (T31.north west);
\draw[dag_arrow] (T31.south) -- (TB1.north);
\draw[dag_arrow] (C2C8.south) to[out=-90,in=110,looseness=1.05] (TB1.north west);
\draw[dag_arrow] (C11.south) -- (Examples.north);
\draw[dag_arrow] (TB1.east) -- (Examples.west);

\end{dagPaperBody}
\end{tikzpicture}
\end{document}
